(1) 在圆偏振基 \(\mathbf e_\sigma=(\mathbf e_x+i\sigma\mathbf e_y)/\sqrt2\) 中，局部线性延迟器可写为圆偏振转换加单位项；半波 \(\delta=\pi\) 时输入旋向 \(\sigma=\pm1\) 变成 \(-\sigma\)，同时获得Pancharatnam–Berry几何相位 \(e^{2i\sigma\alpha(x)}\)（整体固定相位无物理影响）。

(2) 因为 \(\alpha(x)=qx\)，输出复幅的横向相位为 \(\Phi_\sigma(x)=2\sigma qx\)，故额外横向波数 \(\Delta k_x=\partial_x\Phi=2\sigma q\)。出射方向需满足 \(\sin\theta_\sigma=2\sigma q/k_0\)；若 \(|2q|>k_0\)，对应平面波横向波数超过自由空间光锥，不能将它解释为普通传播光线。

(3) 一般延迟器在圆偏振基的输出可分解为同旋 \(\cos(\delta/2)\mathbf e_\sigma\) 与翻旋 \(i\sin(\delta/2)e^{2i\sigma\alpha}\mathbf e_{-\sigma}\)（全局相位约定可变）。两正交分量的功率比例为 \(\cos^2(\delta/2)\) 和 \(\sin^2(\delta/2)\)，其和严格为1。

(4) 偏振自旋转换与空间光轴纹理共同使翻旋分量附加几何相位和横向动量，所需反作用由图案化器件承担。\(q\to0\) 时相位均匀、无光束偏转；\(\delta\to0\) 时翻旋效率趋零，全部光保持入射旋向，不能只因光轴图案仍在就宣称有光束分束。
